\begin{figure}[!ht]
\centering
\begin{tikzpicture}[alt={Four rows of native observation pairs show first level one and reconstructed DC 23; first inverse output three and sample 131; next level zero and reconstructed DC zero; next inverse output zero and sample 131. The two retained sample values have outlines.}]
\node[anchor=north west,inner sep=0,font=\small\bfseries] (title) at (0,0) {Follow selected levels into the next reference};
\node[anchor=north west,inner sep=0] (a) at ([yshift=-2mm]title.south west) {\MARSRevisionCrop{revision/captures/h263-prediction/h263-chain.png}{0.99\linewidth}{0.03125}{0.26888889}{0.58203125}{0.60888889}};
\begin{scope}[shift={(a.south west)},x={($(a.south east)-(a.south west)$)},y={($(a.north west)-(a.south west)$)}]
\draw[black,line width=.7pt] (0.582500,0.386818) rectangle (0.965000,0.569727);
\draw[black,dashed,line width=.7pt] (0.582500,0.005000) rectangle (0.965000,0.187909);
\end{scope}
\end{tikzpicture}
\caption{The caller copies four pairs from its two actual block results: selected DC and reconstructed DC, then inverse output and sample, first for raw 130 and then for raw 132. The outlined and dashed sample fields are both 131. The second call uses the first result as its prediction.}
\label{fig:revision-h263-native-chain}
\end{figure}
